%!TEX program = lualatex
% SPDX-License-Identifier: CC-BY-SA-4.0
% latex-atlas 04: theorem apparatus after the Principia. Every statement is a
% Proposition with a running number and a kind ("Proposition II. Lemma I."),
% heads are centred spaced capitals, a demonstration closes with Q.E.D., and
% remarks are scholia. Against it, amsthm's plain style. The mathematics is
% set in a math font cut to match the text, against Latin Modern Math.
\documentclass[11pt]{article}

\usepackage{iftex}
\ifLuaTeX\else
  \PackageError{latex-atlas}{This example requires LuaLaTeX}
    {Run lualatex, not pdflatex.}
\fi

% --- lab controls ---------------------------------------------------------
% \atlascontrol{<name>}{<treatment>}{<modern default>} sets the treatment
% unless <name> is in \atlasomit, a comma list build.sh --ablate passes on
% the command line. \atlasstamp records the controls at the foot of the plate
% so every variant names itself. Identical in every example, by design.
\ExplSyntaxOn
\providecommand{\atlasomit}{}
\clist_new:N \l_atlas_omit_clist
\clist_new:N \g_atlas_seen_clist
\clist_set:Ne \l_atlas_omit_clist { \atlasomit }
\NewDocumentCommand{\atlascontrol}{ m +m +m }
  {
    \clist_gput_right:Nn \g_atlas_seen_clist { #1 }
    \clist_if_in:NnTF \l_atlas_omit_clist { #1 } { #3 } { #2 }
  }
\NewDocumentCommand{\atlasstamp}{}
  {
    \clist_remove_duplicates:N \g_atlas_seen_clist
    Controls:~\clist_use:Nn \g_atlas_seen_clist { ,~ }.
    \clist_if_empty:NF \l_atlas_omit_clist
      { ~This~impression~sets~\clist_use:Nn \l_atlas_omit_clist { ,~ }~to~the~modern~default. }
  }
\ExplSyntaxOff
% --------------------------------------------------------------------------

\usepackage{amsthm}
\usepackage{fontspec}
\usepackage{unicode-math}
\usepackage{xcolor}
\usepackage{geometry}
\usepackage[tracking=true,protrusion=false,expansion=false]{microtype}

\defaultfontfeatures{Renderer=HarfBuzz}
\IfFontExistsTF{EB Garamond}{
  \setmainfont{EB Garamond}[
    Numbers={OldStyle,Proportional},
    Ligatures={TeX,Common,Rare}
  ]
}{
  \setmainfont{TeX Gyre Pagella}[
    Numbers={OldStyle,Proportional},
    Ligatures={TeX,Common}
  ]
}

% Mathematics in a face cut for the text (Garamond-Math beside EB Garamond,
% Pagella Math beside Pagella), or the engine's default, Latin Modern Math.
% Math digits are lining in either case, so the old-style figures of the
% text never leak into a formula.
\atlascontrol{mathfont}{
  \IfFontExistsTF{Garamond-Math.otf}{
    \setmathfont{Garamond-Math.otf}[StylisticSet={5,7}]
  }{
    \setmathfont{TeX Gyre Pagella Math}
  }
}{
  \setmathfont{Latin Modern Math}
}

% Paper tint is per-example tuning for this proposition page, not a shared
% token yet; the examples deliberately differ slightly (see README).
\atlascontrol{paper}{
  \definecolor{parchment}{HTML}{F5ECDB}
  \definecolor{ink}{HTML}{241F19}
}{
  \definecolor{parchment}{HTML}{FFFFFF}
  \definecolor{ink}{HTML}{000000}
}
\pagecolor{parchment}
\color{ink}
\geometry{a4paper,top=30mm,bottom=30mm,left=34mm,right=34mm}
\setlength{\parindent}{1.1em}
\setlength{\parskip}{0pt}
\pagestyle{empty}

% Proposition heads: one running count across kinds, one count per kind,
% both in roman numerals, centred in spaced small capitals; the statement
% in italic; an attribution, if given, closes the statement. The
% demonstration is headed as such. The modern default is amsthm's plain
% style with a single shared counter.
\atlascontrol{heads}{
  \newcounter{proposition}
  \newcounter{theoremkind}
  \newcounter{lemmakind}
  \newcounter{problemkind}
  \NewDocumentCommand{\prophead}{o m}{%
    \par\addvspace{1.1em plus .3em}%
    \centerline{\small\textls[120]{\textsc{#2}}}%
    \IfValueT{#1}{\nobreak\vspace{.1em}\centerline{\footnotesize\itshape #1}}%
    \nobreak\vspace{.3em}\noindent\ignorespaces}
  \NewDocumentEnvironment{Definition}{}
    {\prophead{Definitions.}\upshape\ignorespaces}{\par}
  \NewDocumentEnvironment{Theorem}{o}
    {\refstepcounter{proposition}\refstepcounter{theoremkind}%
     \prophead[#1]{Proposition \Roman{proposition}.\quad Theorem \Roman{theoremkind}.}\itshape\ignorespaces}
    {\par}
  \NewDocumentEnvironment{Lemma}{o}
    {\refstepcounter{proposition}\refstepcounter{lemmakind}%
     \prophead[#1]{Proposition \Roman{proposition}.\quad Lemma \Roman{lemmakind}.}\itshape\ignorespaces}
    {\par}
  \NewDocumentEnvironment{Problem}{o}
    {\refstepcounter{proposition}\refstepcounter{problemkind}%
     \prophead[#1]{Proposition \Roman{proposition}.\quad Problem \Roman{problemkind}.}\itshape\ignorespaces}
    {\par}
  \NewDocumentEnvironment{Corollary}{}
    {\par\addvspace{.5em}\noindent\textsc{Cor.}\ \ignorespaces}{\par}
  \renewcommand{\proofname}{Demonstration}
}{
  \theoremstyle{plain}
  \newtheorem{Theorem}{Theorem}
  \newtheorem{Lemma}[Theorem]{Lemma}
  \newtheorem{Problem}[Theorem]{Problem}
  \newtheorem{Corollary}[Theorem]{Corollary}
  \theoremstyle{definition}
  \newtheorem*{Definition}{Definitions}
}

% The close of a demonstration: Q.E.D. in spaced small capitals, or the
% tombstone.
\atlascontrol{qed}{
  \renewcommand{\qedsymbol}{\textls[60]{\textsc{q.e.d.}}}
}{}

% Scholium: a remark set smaller, headed in spaced small capitals; the
% modern default is amsthm's Remark.
\atlascontrol{scholium}{
  \NewDocumentEnvironment{Scholium}{}
    {\par\addvspace{.7em}\small\noindent\textls[100]{\textsc{Scholium.}}\ \ignorespaces}{\par}
}{
  \theoremstyle{remark}
  \newtheorem*{Scholium}{Remark}
}

\begin{document}

\begin{center}
  {\small\textls[130]{\textsc{Of the Pursuit upon a Graph}}}\par
  \vspace{0.3em}
  {\Large\itshape Cops and Robber, in Propositions}\par
  \vspace{0.3em}
  {\footnotesize after the manner of the \textit{Principia}}
\end{center}

\vspace{0.6em}

\begin{Definition}
I. A vertex $u$ of a graph $G$ is a \emph{corner} when some other vertex
$v$ sees all that $u$ sees, that is, $N[u] \subseteq N[v]$.
II. $G$ is \emph{dismantlable} when its vertices may be struck out one at
a time, each a corner of the graph that remains, until a single vertex is
left.
III. The \emph{cop number} $c(G)$ is the least number of cops which,
moving by turns with the robber along the edges or resting, are certain
of his capture; and $G$ is \emph{cop-win} when $c(G) = 1$.
\end{Definition}

\begin{Theorem}[Nowakowski and Winkler; Quilliot, 1983]
A finite connected graph is cop-win if and only if it is dismantlable.
\end{Theorem}

\begin{proof}
Let $u$ be a corner dominated by $v$, and let $G' = G - u$. The map that
sends $u$ to $v$ and fixes the rest is a retraction of $G$ upon $G'$. If
the cop wins upon $G'$, let him play there against the robber's shadow,
which is the robber's vertex with $u$ read as $v$; when the shadow is
taken the robber stands at the cop's vertex or at $u$, and $v$ commands
both. Conversely, upon a cop-win graph the robber's last vertex before
capture is dominated by the cop's; strike it out, and the graph that
remains is cop-win again, whence the induction.
\end{proof}

\begin{Corollary}
Every finite tree is cop-win, a leaf being a corner; and so is every
chordal graph, a simplicial vertex being one likewise.
\end{Corollary}

\begin{Lemma}[Berarducci and Intrigila, 1993]
If $H$ is a retract of $G$, then $c(H) \le c(G)$.
\end{Lemma}

\begin{proof}
Let $r\colon G \to H$ be the retraction, and let $c(G)$ cops win upon
$G$. Upon $H$ let them play the same strategy against the robber, whose
vertex, lying already in $H$, is its own image under $r$; every move the
strategy prescribes lies in $G$, and is carried into $H$ by $r$, which
preserves adjacency or identity. The capture in $G$ is therefore a
capture in $H$.
\end{proof}

\begin{Problem}[Meyniel, 1985]
Whether there is an absolute constant $C$ such that, for every connected
graph $G$ upon $n$ vertices,
\[
  c(G) \le C\sqrt{n}.
\]
\end{Problem}

\begin{Scholium}
The best bound yet demonstrated, by Lu and Peng, by Scott and Sudakov,
and by Frieze, Krivelevich and Loh, is
$c(G) \le n\,2^{-(1-o(1))\sqrt{\log_2 n}}$, which exceeds every
$n^{1-\varepsilon}$. The conjecture stands; it is settled for graphs of
diameter two, and for random graphs.
\end{Scholium}

\vspace{1.4em}
\footnotesize\noindent
The apparatus of the \textit{Principia}: propositions numbered in one
sequence and named by kind, demonstrations rather than proofs, corollaries
run in, scholia set small. The mathematics is cut to match the text face;
its digits are lining while the text's are old-style, and the two never
meet in one line.

\medskip\noindent\itshape\atlasstamp

\end{document}
